\begin{savequote}
\begin{verbatim}
( Excellent time to become a missing )
( person.                            )
 ------------------------------------
  o
   o
       ___  
     {~._.~}
      ( Y )
     ()~*~()   
     (_)-(_) 
\end{verbatim}
\end{savequote}

\chapter
[\CO{[TOX]} \emph{Spezielle Patientengruppe:} Vergiftungen]
{\CC{[TOX]}\\\emph{Spezielle Patientengruppe:} Vergiftungen}
\minitoc
% %%%%%%%%%%%%%%%%%%%%%%%%%%%%%%%%%%%%%%%%%%%%%%%%%%%%%%%%%%%
% \section{Vergiftungen -- Toxikologie}
\label{topic:vergiftung}

\dictum[Hohenheim von Theoprast]{Die Dosis macht das Gift.}

\Definition{\T{Vergiftung} bedeutet eine (meist  dosisabhängige)
organische  Funktionsbeeinträchtigung bzw. Schädigung des
Organismus durch eine aufgenommene Substanz.}

Vergiftungen sind eine interessante Sache: Ein an sich
nützlicher Stoff kann sich in eine todbringende Waffe
verwandeln, wenn er im Übermaß zugeführt wird. So ist zum
Beispiel ein Überleben ohne Salz nicht möglich. Führt man
jedoch 20\,dag zu, endet das meist tödlich. Sonst nützliche
Medikamente können schon bei einer Dosisabweichung im
Mikrogramm-Bereich eine fatale Wirkung haben.

Andererseits können Vergiftung in Abhängigkeit vom ursächlichen Stoff so
ziemlich alle Erscheinungsformen darbieten. Oft findet man
nur \EE{unspezifische Symptome} vor, wie z.\,B. Erbrechen,
Durchfall, Schmerzen im Magen-Darm-Trakt,
Bewusstseinsstörungen, Krampfanfälle oder Entgleisung der
Vitalparameter.

Es ist weder möglich, noch sinnvoll, auf sämtliche mögliche
Vergiftungen im Detail einzugehen. Im Folgenden werden
einige häufig anzutreffende Vergiftungen sowie
Herangehensweisen für andere, weniger häufige, vorgestellt
und besprochen. 





\Layout{2}{Der Weg in den Körper -- Aufnahme des Giftes}{%
Die Aufnahme der Substanz kann über verschiedene Wege
erfolgen:

\begin{description}
  \item[Oral] über den Mund
  \item[Inhalativ] Einatmen
  \item[Stechen, spritzen] in die Vene (intravenös, z.\,B.
  Drogenkonsum), in den Muskel, in die %(Unter-)
  Haut (z.\,B. Insektenstich)
  \item[Über die Haut] transkutan : fettlösliche Substanz
  kann  Hautbarriere durchdringen (z.\,B.:  organische
  Lösungsmittel,  Insektenschutzmittel etc.)
\end{description}
}{}{}{}{}





\subparagraph*{Vorgehensweise}

{\mdseries
5-Fingerregel:}

\begin{enumerate}
  \item Vitalparameter
  \item Giftentfernung
  \item Antidot-Therapie (Notarzt)
  \item Gift-Asservierung
  \item Transport
\end{enumerate}

\Layout{0}{Vergiftungsinformationszentrale}{%
Bei der Vergiftungsinformationszentrale kann Expertenrat
eingeholt werden, sofern der verantwortliche Stoff bekannt
ist. Unter Umständen kann wird hier auch bei der
Identifizierung des Stoffes geholfen.
\Fazit{\EE{{\ttfamily 01 / 406 43 43}}}
}{
Expertenrat

01 / 406 43 43
}{}{}{}






Gegen die meisten Vergiftungen kann man präklinisch,
ursächlich wenig unternehmen. Das Hauptaugenmerk liegt
daher auf: 

\MassnahmenWide{Allgemeine Maßnahmen und
Herangehensweise}{Vergiftungen}{\label{m:am-vergiftung}}{
\begin{enumerate}
  \item Selbstschutz
  \item Beendigung des Giftkontaktes
  \item Evtl. Dekontamination bei
  Gefahrstoffen\footnote{z.\,B. Insektizide, Lösungsmittel
  in der chemischen Industrie, \ldots}
  \item Sicherung der Vitalfunktionen
  \item Entscheidung, ob es sich um einen akut vital
  bedrohten Patienten handelt
  \item Quellenforschung (Mistkübel
  durchsuchen)
  \item Nüchtern lassen
  \item NICHT absichtlich Erbrechen
  herbeiführen
  \item Asservierung (Packungen, Erbrochenes in
  Nierentasse mitnehmen!)
  \item Anamnese, Besonderheiten:
  \begin{multicols}{2}
\begin{enumerate}\raggedbottom
  \item \EE{Wer}? (Alter, KG, Vorerkrankungen, \ldots
  $\rightarrow$ Risikoabschätzung!)
  \item \EE{Was}? (Art des Giftes $\rightarrow$ ev.
  zu erwartende Symptome)
  \item \EE{Wann}? (Zeit seit Einnahme)
  \item \EE{Wie}?
  (inhalativ, oral, perkutan, intravenös etc.)
  \item \EE{Wieviel}? (Dosisabschätzung)
  \item \EE{Warum}? (Ursachenforschung:
  Selbst\-mord\-ver\-such, Unfall, Fremdverschulden,
  Vorsatz?)
\end{enumerate}
\end{multicols}
\item ggfs. \EE{Expertenrat einholen} und Entscheidung
bezüglich des weiteren Vorgehens

\begin{multicols}{2}
\begin{enumerate}\raggedbottom
  \item Gegenmittel verfügbar?
  \item 'normales' Spitalsbett?
  \item Intensivüberwachung?
  \item Intensivüberwachung mit Entgiftungseinheit?
  \item Dialyse?
  \item Druckkammer?
  \item Sonstige Spezialbehandlungseinheit?
\end{enumerate}
\end{multicols}
\end{enumerate}
}



\section{Vergiftungen mit Alkohol}

\Layout{20}{Einleitung}{%
Die Vergiftung mit Alkohol ist häufig anzutreffen, weil "`sozial akzeptiert"'.
Wenn der Betroffene "`über's Ziel hinaus schießt"', kann es schnell zu
lebensbedrohlichen Zuständen kommen.
}
{}
{./images/TAE/dif-w2-zelt-1.jpg}
{}
{}


\Z{Die akute Alkoholvergiftung ist von der chronischen
Alkoholabhängigkeit (Alkoholkrankheit) zu unterscheiden.}

\Layout{0}{Wirkung}
{
Alkohol  wirkt enthemmend, z.T. stimulierend und
stimmungsmodulatorisch. Weiters kommt es zu einer Sprach- 
und Gangunsicherheit. 
}{
\begin{itemize}
    \item Enthemmend, stimmungsmodulatorisch
    \item Sprach- und Gangunsicherheit
    \item Bewusstseinstrübung
\end{itemize}
}{}{}{}

\Layout{0}{\TxGefahren}
{
Mit zunehmender Menge kommt es zu einem totalen
Kontrollverlust, Bewusstseinstrübung,
Schmerzunempfindlichkeit, Temperaturregulationsstörung,
Inkontinenz\footnote{\emph{Inkontinenz:} Unfähigkeit Harn
und/oder Stuhl zu halten.} Im Extremfall kann es zu einer 
vegetativen Entgleisung und einem tief komatösen Zustand 
kommen. Spätestens dann besteht akute Lebensgefahr. Weitere 
Gefahren sind Atemdepression, Aspiration, Ausfall der 
Schutzreflexe und -- besonders in der kalten Jahreszeit --
Unterkühlung.
}{
\begin{itemize}
    \item Koma
    \item Atemdepression
    \item  Ausfall der Schutzreflexe
    \item Aspiration
    \item Schmerzunempfindlichkeit
    \item Unterkühlung
\end{itemize}
}{}{}{}


Es ist unbedingt zu erheben, ob noch andere Substanzen
(Medikamente, Drogen, \ldots) eingenommen wurden
(\E{"`Mischintoxikation"'}). Aufgrund der oft reduzierten
Schmerzempfindlichkeit ist besonders auf Verletzungen zu
achten.

Die \EE{vitale Bedrohung} ist vor allem anhand der
Vitalparameter, der körperlichen Untersuchung und des
Bewusstseinszustandes zu beurteilen.

\MassnahmenWide{Spezielle Maßnahmen}{Vergiftungen mit Alkohol}{}{
\begin{itemize}
    \item Allgemeine Maßnahmen bei Vergiftungen
    (\prettyref{m:am-vergiftung})
    \item \TxBeiVit \TxMassVit
\end{itemize}
}


Der Alkoholentzug und das Entzugsdelir wird unter
\prettyref{topic:alkohol-entzug} abgehandelt.

\Layout{2}{Chronische Schäden}
{
\begin{itemize}
    \item Leberschäden, Zirrhose
    \item Aszites (Wasserbauch)
    \item Krampfadern in der unteren Speiseröhre  (Blutungsgefahr)
    \item Verminderte Blutgerinnung
    \item Vitaminmangel
    \item Geistiger Abbau, Gedächtnisschwäche
    \Z{(Korsakow-Syndrom, Wernicke-Enzephalopathie)}
    \item Anämie
    \item Pankreatitis
    \item Fetales Alkoholsyndrom (Schädigung des Fötus in der Schwangerschaft)
    \cite{Jones1973,Guerri2009}
\end{itemize}
}{	

}{}{}{}













\section{Vergiftungen durch Medikamente}

\Layout{60}{}
{%
Überdosierungen von Medikamenten können komplett
unterschiedliche Symptome erzeugen, je nach eingenommenen
Medikament. Manchmal können die Symptome sogar erst Stunden
nach der Einnahme auftreten und tödlich enden. Vergiftungen
mit Medikamenten gehören daher immer abgeklärt und die
Patienten hospitalisiert. 

Neben unabsichtlichen Fehleinnahmen oder Überdosierungen
werden Arzneimittel oft auch in Selbstmordabsicht
eingenommen. Die Umstände der Einnahme müssen daher
abgeklärt werden.
}{

}{./images/teaser/imgp0482_00800.jpg}{}{Gabriel}

\Z{
\paragraph{Paracetamol (Mexalen\texttrademark,
APA\texttrademark, \ldots)} Paracetamol kann in hoher Dosis
zu Leberversagen führen. Symptome treten erst nach Stunden
auf. Es gibt ein Gegenmittel, dieses muss jedoch
rechtzeitig gegeben werden. Ohne Behandlung kommt es zum
Leberversagen und in Folge zum Tod. 

\paragraph{Digitalis} Digitalis ist ein auf das Herz
wirkendes Medikament und wird seit langer Zeit zur
Behandlung der Herzinsuffizienz und von
Herzrhythmusstörungen eingesetzt. Es ist jedoch sehr
schwierig die richtige Dosis für einen Patienten zu finden,
daher kommt es -- trotz korrekter Einnahme -- häufig zu
Überdosierungen. Die Folge sind Herzrhythmusstörungen und
andere EKG-Veränderungen und "`Farbensehen"'.
}

\section{Vergiftungen mit Suchtmitteln und 'Drogen'}

\subsection{'Fixer' -- Heroin \& Co.}

Heroin und Morphium gehören zu der Gruppe der Opiate und
sind im Drogenmilieu weit verbreitet, die Anwendung erfolgt
i.\,d.\,R. intravenös. 

\Layout{0}{\TxSymptome{} Opiatvergiftung}
{%
Die gefährlichste Wirkung der Opiate ist die
Atemdepressionen, d.h. die Atemfrequenz wird dosisabhängig
immer weiter reduziert bis es schließlich zu einem
Atemstillstand kommt. Die Patienten zeigen
dementsprechend Zeichen einer Hypoxie
(Zyanose). Charakteristisch ist die starke Verengung der
Pupillen \E{("`stecknadelkopfgroße
Pupillen"')}\index{Pupillen!stecknadelkopfgroße}. Oft
findet man alte Einstichstellen in der Armbeuge. Da in der
"`Szene"' oft unsauberes oder getauschtes Spritzbesteck
verwendet wird, besteht eine erhöhte Erkrankungsrate an HIV
und Hepatitis C. Eine weitere Gefahr geht vom Spritzbesteck
aus, wenn es nicht sicher verwahrt wird (\Ca{Vorsicht vor
der Nadel!})
}{
\begin{itemize}
    \item Atemdepression
    \item Hypoxie
    \item Stecknadelkopfgroße Pupillen
    \item Oft Begleiterkrankungen (HIV, Hep. C) 
    \item \Ca{Vorsicht Spritzbesteck!}
    \item Gegenmittel verfügbar, wirkt aber nicht so lang
\end{itemize}
}{}{}{}

\Layout{2}{Gegenmittel}
{%
Jedes Notarztmittel führt ein Opiat-Gegenmittel mit sich,
dieses ist sehr effektiv\footnote{\Z{Naloxon. Im Handel
als: \emph{Narcanti}\texttrademark,
\emph{Naloxon-ratiopharm}\texttrademark, u.\,a.}}. Es ist
jedoch wichtig zu wissen, dass dieses Gegenmittel nicht so
lange wirkt wie das Opiat, d. h. es kann wieder zu einer
verminderten Atmung kommen, wenn die Wirkung des Gegenmittels nachlässt. Daher soll der Transport nur mit Arztbegleitung erfolgen.

}{

}{}{}{}


\MassnahmenNG{Spezielle Maßnahmen}{Opiatvergiftung}{}{
\begin{itemize}
    \item Allgemeine Maßnahmen bei Vergiftungen
    (\prettyref{m:am-vergiftung})und
    symptomatische Therapie
    \item \TxBeiVit \TxMassVit
\end{itemize}

%Allgemeine Maßnahmen bei Vergiftungen
%(\ref{vergiftung-allgemein-massnahmen})und symptomatische
%Therapie. 
% Gegenmittel verfügbar.
}

\subsection{'Clubber' -- Ecstacy, Schwammerl,
Special K \ldots}

\Layout{0}{\TxBeschreibung}
{
\Z{MDMA (3,4-Methylen\-dioxy\-meth\-amphet\-amin),} bekannt
unter dem Namen \E{"`Ecstasy"' (XTC, E, X)}, ist eine der
bekanntesten \E{"`Partydrogen"'}. Auf den ersten Blick
scheint die Substanz auch perfekt für diesen Zweck zu sein:
Der Patient erlebt ein gesteigertes Glücksgefühl, ein
Gefühl von unerschöpflicher Energie und ein gesteigertes
Erleben der Umwelt. 
}{
\begin{itemize}
  \item Ecstasy"' (XTC, E, X).
  \item Gesteigertes "`Glücksgefühl"'.
  \item Gefühl von unerschöpflicher Energie.
  \item Gesteigertes Erleben der Umwelt.
\end{itemize}
}{}{}{}

\Layout{0}{\TxSymptome}
{%
Die Nebenwirkungen sind jedoch
gravierend und können lebensgefährlich sein. Man kann sie
unterteilen in körperliche und psychische Nebenwirkungen. 
Psychisch kommt es dosisabhängig zu Halluzinationen
und Panik-Attacken. 

Körperlich können Patienten tachykarde
Herzrhythmusstörungen zeigen, Krampfanfälle, hypertensive
Krisen, massiver Körpertemperaturanstieg über
40\textdegree{} und Muskelzerfall.

Bei einer ausgeprägten Vergiftung kann eine vitale
Bedrohung bestehen. 

\vspace{1cm}

}{
\begin{itemize}
  \item Nebenwirkungen Psychisch
  \begin{itemize}
    \item Halluzinationen
    \item Panik-Attacken
  \end{itemize}
  \item Nebenwirkungen Körperlich
  \begin{itemize}
    \item Tachykarde Herzrhythmusstörungen
    \item Hypertensive Krisen
    \item Krampfanfälle
    \item Massiver Temperaturanstieg über
    {\textgreater}\,40\,\textcelsius{}
  \end{itemize}
\end{itemize}
}{}{}{}


\MassnahmenNG{Spezielle Maßnahmen}{Vergiftung mit
Szenedrogen}{}{
\begin{itemize}
    \item Allgemeine Maßnahmen bei Vergiftungen
    (\prettyref{m:am-vergiftung}) und
    symptomatische Therapie
    \item \TxBeiVit \TxMassVit
\end{itemize}

%Allgemeine Maßnahmen bei Vergiftungen
%(\ref{vergiftung-allgemein-massnahmen})und symptomatische Therapie 
}


%„Liquid Ecstasy”: Gamma-Butyrolacton-Entzugsdelir mit
%Rhabdomyolyse und dialysepflichtiger Niereninsuffizienz
\cite{Supady2009}









\section{Gase}

\subsection{Vergiftungen mit Stickgasen}

\MassnahmenNG{Allgemeine Maßnahmen}{Stickgasvergiftungen}{\label{m:am-stickgase}}{


\begin{enumerate}
  \item Auf Selbsschutz achten!
  \item Patienten retten, aber nur wenn es der
  Eigenschutz zulässt.
  \item \ce{O2}-Berieselung mit Maske ist in keinem Fall
  ein ausreichender Atemschutz!
  
  \item Spezialkräfte anfordern (Schwerer Atemschutz). I.d.R.
  Feuerwehr
  \item Weitere Opfer ausschließen (lassen):
  Nachbarwohnungen, etc.
  \item Sicherung der Vitalfunktionen, Beurteilung ob der
  Patient vital bedroht ist.
  \item \TxBeiVit \TxMassVitBes
  \begin{itemize}
    \item \ce{O2} hochdosiert, 15\LMin
    \item Lagerung: situationsgerecht.
  \end{itemize}
  \item Traumacheck (Folgeverletzungen)
  \item Neurocheck inkl. Blutzuckerbestimmung
  (Differentialdiagnosen, Beurteilung der Beeinträchtigung im Verlauf)
\end{enumerate}
}


\subsubsection{Kohlenmonoxid (\ce{CO}) -- oder: Von defekten
Heizlüftern, Gasthermen und Autoabgasen}

\Layout{60}{Einleitung}{%
Kohlenmonoxid (\ce{CO}) entsteht bei unvollständiger
Verbrennung (\ce{O2}-Zufuhr ${\downarrow}$), dies passiert
vor allem bei schlecht gewarteten Gasheizungen und
Durchlauferhitzern. \ce{CO} bindet sich 300 mal besser an den
roten Blutfarbstoff Hämoglobin als $O2$, daher verdrängen
schon kleine Mengen den Sauerstoff aus dem Blut! Die Folge
ist logischerweise eine Unterversorgung mit Sauerstoff.
}
{}
{./images/gabriel-sebastian-aass/Atemschutz_IMG_2043_00800.jpg}
{}
{GABS}




\Layout{2}{Anamnese und Szene}
{
\begin{itemize}
    \item Mehrere, unerklärliche Opfer, leblose Haustiere
    \item Gasofen, Durchlauferhitzer
    \item \emph{"`Immer beim duschen!"'}
\end{itemize}
}{

}{}{}{}

\Layout{0}{\TxSymptome}
{% 
Kohlenmonoxid-beladenes Blut ist -- wie
Sauerstoff-beladenes -- hellrot. Der Patient ist daher eher
nicht bleich oder zyanotisch, sondern die Haut ist oft rosig.
Pulsoxymeter älterer Bauart verwechseln
Kohlenmonoxid-beladenes und Sauerstoff-beladenes Blut, daher
wird ein falsch hoher Sauerstoffsättigungswert angezeigt!

Es stehen die Symptome des Sauerstoffmangels im Vordergrund
(Hypoxie): Der Patient leidet Atemnot, allerdings ohne
Zyanose, begleitet von Allgemeinsymptomen wie Schwindel,
Kopfschmerz, Übelkeit und Erbrechen, evtl. besteht auch eine
Tachykardie. 

Es können sich rasch Bewusstseinsstörungen bis hin zur
Bewusstlosigkeit entwickeln, auch Krämpfe sind gut möglich.
}{
\begin{itemize}
  \item Dyspnoe  \E{ohne}  Zyanose
  \item Tachykardie
  \item Schwindel, Kopfschmerz, Übelkeit, Erbrechen
  \item Bewusstseinsstörungen, Koma
  \item Krämpfe
  \item Pulsoxymeter zeigt falsche Werte an!
\end{itemize}
}{}{}{}


\MassnahmenNG{Spezielle Maßnahmen}{Kohlenmonoxid-Vergiftung}{}{

\begin{itemize}
    \item Allgemeine Maßnahmen bei
Stickgas-Vergiftungen
(\prettyref{m:am-stickgase})
\item \Z{Hyperbare $O2$-Therapie in der Druckkammer erwägen}
\end{itemize}
}


% \begin{figure}[H]
%     \begin{WidePar}
% \Parallel{
% \includegraphics[width=\linewidth]{./images/2006-01-31-gaswienheute2.png}
% }{
% \Captionof{figure}{\AuthNotesPicLegalOk{Gas Wien Heute
% Homepage}{OK, wenn richtig titriert}\ce{CO}-Vergiftung.
% Regelmäßig kommt es zu Unfällen durch defekte Gasthermen
% oder Heizlüfter. Bei diesem Einsatz war ein RTW des ASB
% Floridsdorf-Donaustadt beteiligt. Die dritte Patientin
% wurde erst \EE{nachdem die Feuerwehr \E{sämtliche} (!)
% Wohnungen des Wohnhauses aufgebrochen hatte} in der
% \E{darüberliegenden Wohnung} gefunden.
% \SourcePicture{Wien heute \url{http://wien.orf.at}}}
% }
% \end{WidePar}
% 
% \end{figure}

\begin{figure}[H]
\begin{WidePar}
\begin{minipage}{0.6\linewidth}
\hspace{0.1\linewidth}
\includegraphics[width=\linewidth]{./images/2006-01-31-gaswienheute2.png}
\end{minipage}
\begin{minipage}{0.3\linewidth}
\Captionof{figure}{\AuthNotesPicLegalOk{Gas Wien Heute
Homepage}{OK, wenn richtig titriert}\ce{CO}-Vergiftung.
Regelmäßig kommt es zu Unfällen durch defekte Gasthermen
oder Heizlüfter. Bei diesem Einsatz war ein RTW des ASB
Floridsdorf-Donaustadt beteiligt. Die dritte Patientin
wurde erst \EE{nachdem die Feuerwehr \E{sämtliche} (!)
Wohnungen des Wohnhauses aufgebrochen hatte} in der
\E{darüberliegenden Wohnung} gefunden.
\SourcePicture{Wien heute \url{http://wien.orf.at}}}
\end{minipage}
\end{WidePar}
\end{figure}

\subsubsection{Kohlendioxid (\ce{CO2}) -- Tod im Weinkeller}


Kohlendioxid fällt als Gär- oder Verbrennungsgas gehäuft in
der Landwirtschaft und der Industrie an. Es ist schwerer als
Luft und sammelt sich in Mulden, Weinkellern, Silos, etc.


\Layout{2}{\TxSymptome}
{%
\begin{itemize}
  \item Kopfschmerz, Schwindel, Übelkeit
  \item Krämpfe
  \item Zyanose
  \item Hyperventilation, da der hohe
  \ce{CO2}-Spiegel im Blut das Atemzentrum stimuliert
\end{itemize}
}{

}{}{}{}


\MassnahmenNG{Spezielle Maßnahmen}{Kohlendioxid-Vergiftung}{}{
Allgemeine Maßnahmen bei
Stickgas-Vergiftungen
(\prettyref{m:am-stickgase})
}

\Cave{
Bei Gasunfällen hat der \EE{Selbstschutz} Vorrang.
Oft ist die Bergung nur mit \EE{schwerem
Atemschutz} möglich ${\rightarrow}$ FEUERWEHR!}

\Fazit{Ein toter Helfer ist keine Hilfe.}

\subsection{Reizgase}

\Layout{2}{Typen}
{%
\begin{description}
  \item[Soforttyp] (z.\,B. Tränengas,...) Augenrötung, starker Hustenreiz
  \item[Verzögerungstyp] \Z{(z.\,B. Nitrosegase)} kaum
  Hustenreiz, Latenzphase (kaum Symptome),
  \EE{\VonBis{3}{24}\,h später: toxisches Lungenödem}; bei
  Nichtbehandlung evtl. irreversible Lungenschäden
\end{description}
}{

}{}{}{}

\Layout{46}{Vorkommen}
{

}{

}{
\begin{tabulary}{\linewidth}{LL}
\TabHeading{Vorkommen} & \TabHeading{Art}
\\\midrule
Chemische Industrie 
&
Nitrosegase,  Schwefelwasserstoffe, Ammoniak
\\\addlinespace
Kunststoffindustrie 
&
Chlorwasserstoffe,  Formaldehyd
\\\addlinespace
Düngemittelherstellung 
&
Ammoniak,  Nitrosegase, Schwefelwasserstoffe
\\\addlinespace
Haushalt 
&
Chlorgas (Kombination von best.  Haushaltsreinigern mit
Essig), Lacke,  Imprägniersprays, ...
\\\bottomrule
\end{tabulary}
}{Vorkommen von Reizgasen.}{}





\MassnahmenNG{Allgemeine
Maßnahmen}{Reizgasvergiftungen}{\label{m:am-reizgase}}{


\begin{enumerate}
  \item Auf Selbstschutz achten!
  \begin{itemize}
    \item Patienten retten, aber nur wenn es der
    Eigenschutz zulässt.
    \item \ce{O2}-Berieselung mit Maske ist in keinem Fall
    ein ausreichender Atemschutz!
  \end{itemize}
  \item Gefahr erkennen: Welcher Stoff? Expertenrat einholen!
  \item Wenn erforderlich: Absperrmaßnahmen durchführen
  \item Wenn erforderlich:  Spezialkräfte anfordern
  \item Vorgehen analog zu Abschnitt "`Gefahrengutunfall"'
  (\prettyref{topic:gefahrengutunfall})
  \item Weitere Opfer ausschließen (lassen):
  Nachbarwohnungen, etc
  \item Sicherung der Vitalfunktionen, Beurteilung ob der
  Patient vital bedroht ist.
  \item \TxBeiVit \TxMassVitBes
  \begin{itemize}
    \item \ce{O2} hochdosiert, 15\LMin
    \item Lagerung: situationsgerecht
  \end{itemize}
  \item Traumacheck (Folgeverletzungen)
  \item Neurocheck inkl. Blutzuckerbestimmung
  (Differentialdiagnosen, Beurteilung der Beeinträchtigung im Verlauf)
\end{enumerate}
}




\subsection{Kampfgase}

\Layout{60}{}
{%
Kurz und bündig: Hier gehören Spezialisten her.
Selbstschutz und Expertenrat einholen.
}{
}{./images/gabriel-sebastian-aass/schild-atemschutzmaske-abnehmen.jpg}{}{}


\clearpage
\section{Wovon man nicht trinken sollte \ldots}

\subsection{Einnahme von Säuren und Laugen}
\label{topic:veraetzung-intoxikation}

Verletzungen der Haut und Augen durch Säuren und
Laugen werden Kapitel \emph{"`Traumatologische Notfälle"'}
unter \prettyref{topic:veraetzung-trauma} besprochen.

\Definition{\T{Verätzung}: Lokale Schädigung der
Oberfläche aufgrund einer irreversiblen Zerstörung (Denaturierung) von
Eiweißstoffen.}

% 
% Betroffene Organe
% 
% \begin{itemize}
%     \item Magen-Darm-Trakt
%     \item Haut
%     \item Augen
% \end{itemize}

\Layout{2}{Symptome bei Einnahme}
{
\begin{itemize}
    \item Ätzspuren (Mund/Rachen), Abrinnspuren
    \item Würgen
    \item Erbrechen
    \item Schmerzen
    \item Vorsicht: Ruptur/Perforation/Blutungsgefahr!
\end{itemize}
}{

}{}{}{}




\MassnahmenNG{Spezielle Maßnahmen}{Einnahme von Säuren oder Laugen}{}{
\begin{itemize}
  \item Allgemeine Maßnahmen bei Vergiftungen
  (\prettyref{m:am-vergiftung})
  \item Schluckweise Wasser zur Verdünnung trinken lassen
  \item \EE{Niemals absichtlich Erbrechen  herbeiführen!}
  \item \Z{Keine Neutralisationsversuche! (Milch
  obsolet!)}
\end{itemize}
}


% Verätzung der Haut \prettyref{topic:veraetzung}



\subsection{Schaumbildner (Wasch-/Putzmittel)}


\Layout{2}{\TxBeschreibung}
{%
Schaumbildner vermindern die Oberflächenspannung, dadurch
werden entspannende Schaumbäder mit Tannenwaldduft möglich,
genauso wie eine effektive Geschirrwäsche. }{

}{}{}{}

\Layout{2}{\TxGefahren{} und \TxKomplikationen{}}
{
\begin{itemize}
    \item Verminderte Oberflächenspannung: 
    Aspirationsgefahr ${\uparrow}$.
    \item \EE{Verlegung der Atemwege}.
\end{itemize}
}{

}{}{}{}

\Layout{2}{\TxSymptome}
{
\begin{itemize}
    \item Schaum
    \item Übelkeit/Erbrechen
\end{itemize}
}{

}{}{}{}







\MassnahmenNG{Spezielle Maßnahmen}{Einnahme von
schaumbildenden Substanzen}{}{
\begin{itemize}
    \item Allgemeine Maßnahmen bei Vergiftungen
    (\prettyref{m:am-vergiftung})
    \item Patienten nüchtern lassen
    \item \EE{Niemals absichtlich Erbrechen 
    herbeiführen!}
\end{itemize}
}





% 
% 
% \subsection{\Z{Reserviert}\A{Vergiftungen mit
% Insektiziden}}
% % TODO Vergiftungen mit Insektiziden
% 
% \subsection{\Z{Reserviert}\A{Vergiftungen durch giftige
% Tiere}}
% % TODO Vergiftungen durch giftige Tiere
% 
% 
% 
